\section{Summary Assessment: Part 1 vs. Part 2}

Chapter 7 applied four criteria to both parts. The table below brings the four verdicts together.

\begin{center}
\renewcommand{\arraystretch}{1.3}
\begin{tabular}{>{\raggedright\arraybackslash}p{0.16\linewidth} >{\raggedright\arraybackslash}p{0.36\linewidth} >{\raggedright\arraybackslash}p{0.36\linewidth}}
\toprule
\textbf{Criterion} & \textbf{Part 1} & \textbf{Part 2} \\
\midrule
Specification Satisfaction & Satisfied what was explicitly asked; missed what was assumed, sometimes for the whole phase. & Caught more gaps earlier; treated inherited gaps as work to complete rather than debt to carry. \\
Test Validity & Assessed informally, feature by feature; a passing suite could still check too little or the wrong thing. & A structural requirement for every feature; a passing result tied explicitly to a stated requirement. \\
Architectural Conformance & Checked two of six categories, only once a problem had already spread. & Checks all six by rule, demonstrated in full for at least one feature. \\
Efficiency and Effort (Human vs. Agent) & Human diagnosed problems; agent executed fixes once told. & Agents surfaced more problems themselves; human's role shifted to approving what was already found. \\
\bottomrule
\end{tabular}
\end{center}

Across all four criteria, the same shift recurs: from problems caught late, informally, and mostly by a human, to problems caught earlier, by rule, with the work of detection shared more evenly between human and agent.

